% =============================================================================
%  proof_complete_standalone.tex
%
%  Self-contained version of proof_complete.tex.  It reuses the same preamble
%  (includes.tex / MacrosAndFigures.tex / Bib.bib) as main.tex.
%
%  Layout:
%    1. Background        -- general facts on symmetric powers of torsors
%                            (over an abstract base) that the proof quotes;
%    2. Geometric setting -- specialise to a curve: bounded ramification,
%                            the symmetric tower over U^(d), compactification,
%                            blowup;
%    3. The master diagram;
%    4. The unramifiedness theorem and its proof.
%
%  Compile with:  latexmk -pdf proof_complete_standalone.tex   (uses biber)
% =============================================================================
\documentclass[11pt]{article}
\usepackage{geometry}\geometry{margin=1in}
\usepackage[T1]{fontenc}
\usepackage{ifthen}

\setlength{\parindent}{1em}
\setlength{\parskip}{0.5em}

\input{includes}
\input{MacrosAndFigures}

\usepackage{csquotes}
\usepackage[
backend=biber,
style=alphabetic,
sorting=ynt
]{biblatex}
\addbibresource{Bib.bib}

\newboolean{ThesisIntro}
\setboolean{ThesisIntro}{false}

\title{Unramifiedness of $\cP^{[d]_G}$ at the exceptional divisor\\[2pt]
       {\large (standalone excerpt)}}
\author{Assaf Marzan}
\date{\today}

\begin{document}
\maketitle

\noindent
\textbf{Roadmap.} \Cref{sec:background} collects, without re-deriving them, the
facts about symmetric powers of torsors that the proof uses, stated over an
abstract base.  \Cref{sec:setting} specialises this to a smooth projective
curve: it recalls the notion of bounded ramification and fixes the geometric
data (the symmetric tower over $U^{(d)}$, the compactification $C' \to C$, the
blowups at the points at infinity).  \Cref{sec:master} assembles everything into
a single master diagram, and \Cref{sec:main} states and proves the
unramifiedness theorem.

% #############################################################################
\section{Background}\label{sec:background}
% #############################################################################

Throughout this section $S$ is a base scheme, $G$ a finite abelian group,
$H \subseteq G$ a subgroup, and $d \geq 0$ an integer.  All bases are assumed
Zariski-locally quasi-projective over $S$, so that the symmetric-power quotients
below exist as schemes.

%------------------------------------------------------------------------------
\subsection{Symmetric powers of torsors}\label{subsec:sym-powers}
%------------------------------------------------------------------------------

Let $\cP$ be a $G$-torsor over an $S$-scheme $X$.  For each
$i \in \llbracket 1,d\rrbracket$ write $p_i : X^{\times_S d} \to X$ for the
$i$-th projection and consider the $G$-torsor
\[
p_1^{-1}\cP \otimes^G \cdots \otimes^G p_d^{-1}\cP \;=\; \KG \backslash \cP^{\times_S d}
\]
over $X^{\times_S d}$, where $\KG \subseteq G^d$ is the kernel of the
multiplication morphism $G^d \to G$ (the dependence on $d$ is suppressed in the
notation).  The symmetric group $S_d$ acts on the right on
$\KG \backslash \cP^{\times_S d}$ by $(p_i)_i \cdot \sigma = (p_{\sigma(i)})_i$,
and this action commutes with the left $G$-action.

\begin{prop}[\cite{Guignard2018}, Proposition 2.32]\label{prop:symmetric_power_torsor}
The right action of $S_d$ on $\KG \backslash \cP^{\times_S d}$ is admissible, so
that the quotient $\cP^{[d]}$ of $\KG \backslash \cP^{\times_S d}$ by $S_d$
exists as a scheme over $S$.  Moreover, the canonical morphism
$\cP^{[d]} \to \mathrm{Sym}_S^d(X)$ is a $G$-torsor, and the morphism
\[
p_1^{-1}\cP \otimes^G \cdots \otimes^G p_d^{-1}\cP \to r^{-1}\cP^{[d]},
\]
where $r : X^{\times_S d} \to \mathrm{Sym}_S^d(X)$ is the canonical projection,
is an isomorphism of $G$-torsors over $X^{\times_S d}$.
\end{prop}

When $G$ is abelian the condition $\prod_{i=1}^d g_i = 1$ defining $\KG$ is
invariant under permuting the coordinates, so the (left) $S_d$-action normalises
the $\KG$-action.  Writing $\Xi = \KG \rtimes S_d$ for the resulting finite
group, acting on $\cP^{\times_S d}$ by
\[
\bigl((g_1,\dots,g_d),\sigma\bigr) \star (p_1,\dots,p_d)
   = \bigl(g_1 p_{\sigma^{-1}(1)},\,\dots,\,g_d p_{\sigma^{-1}(d)}\bigr),
\]
one has $(\KG \backslash \cP^{\times_S d})/S_d \cong \Xi \backslash \cP^{\times_S d}$.

\begin{cor}\label{cor:sym-power-quotient}
There is a natural canonical morphism
\[
q : S_d \backslash \cP^{\times_S d} \longrightarrow \Xi \backslash \cP^{\times_S d}
\qquad\text{equivalently}\qquad
q : S_d \backslash \cP^{\times_S d} \longrightarrow (\KG \backslash \cP^{\times_S d})/S_d ,
\]
the comparison between the symmetric power of $\cP$ \emph{as a scheme} and the
symmetric power of $\cP$ \emph{as a $G$-torsor}.  We write $\cP^{(d)}$ for the
former and $\cP^{[d]}$ for the latter, so that $q : \cP^{(d)} \to \cP^{[d]}$.
\end{cor}
\begin{proof}
$S_d$ is a subgroup of $\Xi$, hence any $\Xi$-invariant morphism is
automatically $S_d$-invariant.
\end{proof}

%------------------------------------------------------------------------------
\subsection{The canonical decomposition of a torsor}\label{subsec:decomposition}
%------------------------------------------------------------------------------

\begin{prop}\label{prop:quotient_torsors_finite}\label{prop:quotient_torsors}
Let $\cP \to X$ be a $G$-torsor in $X_{\text{\'et}}$ and let $H \subseteq G$ be
a subgroup, with quotient map $\pi : G \to G/H$.  Then $\cP \to X$ factors
\[
\cP \xrightarrow{\;\phi\;} \cP_{G/H} \xrightarrow{\;\psi\;} X,
\]
uniquely up to a unique isomorphism commuting with the projections, where
\begin{enumerate}
    \item $\cP_{G/H} := \cP \times^G (G/H)$ is a $G/H$-torsor over $X$, and
    \item $\cP$ is an $H$-torsor over the scheme $\cP_{G/H}$ via $\phi$.
\end{enumerate}
\end{prop}

%------------------------------------------------------------------------------
\subsection{The tower of symmetric quotients}\label{subsection:quotient_torsors_towers}
%------------------------------------------------------------------------------

Fix a $G$-torsor $\cP \to X$ and a subgroup $H \subseteq G$, with canonical
factorisation $\cP \xrightarrow{\phi} \cP_{G/H} \xrightarrow{\psi} X$ as in
\Cref{prop:quotient_torsors_finite} ($\phi$ an $H$-torsor, $\psi$ a
$G/H$-torsor).  We record how the symmetric-power construction of
\Cref{subsec:sym-powers} interacts with this factorisation.

\paragraph{Summation kernels.}
Besides $\KG = \ker(G^d \to G)$ we need
\[
K_H := \ker(H^d \to H), \qquad K_{G/H} := \ker\bigl((G/H)^d \to G/H\bigr).
\]
The snake lemma applied to
\[
\begin{tikzcd}[row sep=large, column sep=large]
0 \arrow[r] & H^d \arrow[r] \arrow[d, twoheadrightarrow] & G^d \arrow[r] \arrow[d, twoheadrightarrow] & (G/H)^d \arrow[r] \arrow[d, twoheadrightarrow] & 0 \\
0 \arrow[r] & H \arrow[r] & G \arrow[r] & G/H \arrow[r] & 0,
\end{tikzcd}
\]
(the vertical maps being the surjective summation morphisms) gives a short exact
sequence
\begin{equation}\label{eq:summation_kernels_ses}
0 \to K_H \to \KG \to K_{G/H} \to 0,
\end{equation}
hence $\KG/K_H \cong K_{G/H}$.  Inside $G^d$ we have the chains of subgroups
\[
K_H \;\subseteq\; \KG \;\subseteq\; H^d + \KG, \qquad H^d \;\subseteq\; H^d + \KG .
\]

\paragraph{Four quotient schemes.}
Let $S_d$ act on $\cP^{\times_S d}$ as in \Cref{cor:sym-power-quotient}, and let
each of the subgroups above act diagonally via the $G$-action on $\cP$.  Set
\[
\begin{aligned}
    \cP^{[d]_G}        &:= (\KG \rtimes S_d)\backslash \cP^{\times_S d},
        &&\text{the symmetric $G$-power of $\cP$;}\\
    \cP^{[d]_H}        &:= (K_H \rtimes S_d)\backslash \cP^{\times_S d},
        &&\text{the symmetric $H$-power of $\cP$ relative to $\phi$;}\\
    (\cP_{G/H})^{(d)}  &:= (H^d \rtimes S_d)\backslash \cP^{\times_S d},
        &&\text{the symmetric power of the scheme $\cP_{G/H}$;}\\
    (\cP_{G/H})^{[d]}  &:= \bigl((H^d + \KG)\rtimes S_d\bigr)\backslash \cP^{\times_S d},
        &&\text{the symmetric $G/H$-power of $\cP_{G/H}$.}
\end{aligned}
\]

\begin{prop}\label{prop:torsor_tower_diagram}
The four quotient schemes above exist and fit into a canonical commutative
diagram
\[
\begin{tikzcd}[row sep=large, column sep=large]
    \cP^{[d]_H} \arrow[r, "H"] \arrow[d] & (\cP_{G/H})^{(d)} \arrow[d, "q"] \\
    \cP^{[d]_G} \arrow[r, "H"'] \arrow[dr, "G"', bend right=20] & (\cP_{G/H})^{[d]} \arrow[d, "G/H"] \\
    & X^{(d)},
\end{tikzcd}
\]
in which each labelled arrow is a torsor under the indicated constant group, the
right-hand vertical $q$ is the canonical projection of
\Cref{cor:sym-power-quotient} applied to $\cP_{G/H} \to X$, and the upper square
is Cartesian:
\[
\cP^{[d]_H} \cong \cP^{[d]_G} \times_{(\cP_{G/H})^{[d]}} (\cP_{G/H})^{(d)}.
\]
The unlabelled left-hand vertical $\cP^{[d]_H} \to \cP^{[d]_G}$ is then the base
change of $q$ along $\cP^{[d]_G} \to (\cP_{G/H})^{[d]}$.  Neither $q$ nor this
left-hand vertical is, in general, a torsor under a constant group scheme.
\end{prop}
\begin{proof}
The four labelled arrows are torsors under the indicated constant groups by
applying \Cref{prop:symmetric_power_torsor} and
\Cref{prop:quotient_torsors_finite}:
\begin{itemize}
    \item to the $G$-torsor $\cP \to X$, giving the $G$-torsor
          $\cP^{[d]_G} \to X^{(d)}$ (the diagonal arrow);
    \item to the $H$-torsor $\phi : \cP \to \cP_{G/H}$ (with $\cP_{G/H}$ as
          base), giving the $H$-torsor $\cP^{[d]_H} \to (\cP_{G/H})^{(d)}$
          (the top arrow);
    \item to the $G/H$-torsor $\psi : \cP_{G/H} \to X$, giving the $G/H$-torsor
          $(\cP_{G/H})^{[d]} \to X^{(d)}$ (the bottom-right vertical) and, via
          \Cref{cor:sym-power-quotient}, the comparison map
          $q : (\cP_{G/H})^{(d)} \to (\cP_{G/H})^{[d]}$.
\end{itemize}
The factorisation $\cP^{[d]_G} \xrightarrow{H} (\cP_{G/H})^{[d]} \xrightarrow{G/H} X^{(d)}$
of the diagonal arrow is the canonical decomposition of
\Cref{prop:quotient_torsors_finite} applied to $\cP^{[d]_G} \to X^{(d)}$ and
$H \subseteq G$: both $\cP^{[d]_G} \times^G (G/H)$ and $(\cP_{G/H})^{[d]}$ are
realised as the admissible quotient
$\bigl((H^d + \KG)\rtimes S_d\bigr)\backslash \cP^{\times_S d}$.

For the Cartesian property, set
$F := \cP^{[d]_G} \times_{(\cP_{G/H})^{[d]}} (\cP_{G/H})^{(d)}$.  Base-changing
the $H$-torsor $\cP^{[d]_G} \to (\cP_{G/H})^{[d]}$ along $q$ exhibits $F$ as an
$H$-torsor over $(\cP_{G/H})^{(d)}$.  The universal property of the fibre
product gives a canonical $H$-equivariant morphism $\cP^{[d]_H} \to F$ over
$(\cP_{G/H})^{(d)}$; any morphism of $H$-torsors is an isomorphism, so
$\cP^{[d]_H} \cong F$ and the upper square is Cartesian.  Via this square the
left vertical $\cP^{[d]_H} \to \cP^{[d]_G}$ is identified with the base change
of $q$ along $\cP^{[d]_G} \to (\cP_{G/H})^{[d]}$.
\end{proof}

% #############################################################################
\section{The geometric setting}\label{sec:setting}
% #############################################################################

%------------------------------------------------------------------------------
\subsection{Bounded ramification along a curve}\label{subsec:bounded-ram}
%------------------------------------------------------------------------------

Let $C$ be a smooth projective geometrically connected curve over a field $k$,
let $\mdls$ be a fixed modulus, and put $U = C \setminus \mdls$.  Let
$\cP \to U$ be a $G$-torsor with $G$ finite abelian.  For a closed point $P$ of
$C$, passing to the completion $\widehat{\Oo}_{C,P}$ gives the complete local
field $\widehat{K}_P$, and restricting $\cP$ to $\Spec(\widehat{K}_P)$ yields a
$G$-torsor over a field, hence a disjoint union of spectra of finite separable
field extensions
\[
\cP|_{\Spec(\widehat{K}_P)} \cong \bigsqcup \Spec(F),
\]
in which the extensions $F/\widehat{K}_P$ are all isomorphic and Galois,
$H = \Gal(F/\widehat{K}_P)$ is the stabiliser of a connected component (a
subgroup of $G$), and there are $[G:H]$ components.  Write $H^r$ for the $r$-th
higher ramification group of $F/\widehat{K}_P$.

\begin{definition}\label{definition:local_ramification_boundness}
The extension $F/\widehat{K}_P$ \textit{has ramification bounded by $r$} if the
ramification group $H^r$ is trivial.  We say \textit{$\cP$ has ramification
bounded by $r$ at $P$} if the corresponding local extensions $F/\widehat{K}_P$
satisfy this condition.
\end{definition}

%------------------------------------------------------------------------------
\subsection{The symmetric tower over \texorpdfstring{$U^{(d)}$}{the symmetric power}}
\label{subsec:tower-over-U}
%------------------------------------------------------------------------------

For the rest of the note we fix, in addition to the data above:
\begin{itemize}
  \item a normal subgroup $H \trianglelefteq G$;
  \item a point $p_\infty \in \mdls$ and an integer $d \geq 1$;
\end{itemize}
and we assume that $\cP \to U$ has ramification bounded by $d$ at $p_\infty$ in
the sense of \Cref{definition:local_ramification_boundness} (i.e.\ the $d$-th
higher ramification group of the extension of DVRs at $p_\infty$ is trivial).

Apply the construction of \Cref{subsection:quotient_torsors_towers} with
$S = \Spec k$ and base $X = U$.  This produces the four quotient schemes
\[
\cP^{[d]_H},\qquad \cP^{[d]_G},\qquad (\cP_{G/H})^{(d)},\qquad (\cP_{G/H})^{[d]},
\]
and, by \Cref{prop:torsor_tower_diagram}, the canonical commutative diagram with
Cartesian upper square
\[
\begin{tikzcd}[row sep=large, column sep=large]
\cP^{[d]_H} \ar[r,"H"] \ar[d]
  & (\cP_{G/H})^{(d)} \ar[d,"q"] \\
\cP^{[d]_G} \ar[r,"H"'] \ar[dr,"G"',bend right=15]
  & (\cP_{G/H})^{[d]} \ar[d,"G/H"] \\
  & U^{(d)},
\end{tikzcd}
\]
in which the arrows labelled $H$, $G/H$ and $G$ are torsors under the indicated
finite constant groups (in particular, \'etale).

%------------------------------------------------------------------------------
\subsection{Compactification and blowup}\label{subsec:compactify-blowup}
%------------------------------------------------------------------------------

The $G/H$-torsor $\cP_{G/H} \to U$ compactifies to a finite morphism of smooth
projective curves $f \colon C' \to C$; fix $p'_\infty \in f^{-1}(p_\infty)$.
Taking $d$-fold symmetric powers gives the finite morphism
\[
f^{(d)} \colon C'^{(d)} \to C^{(d)}, \qquad d \cdot p'_\infty \longmapsto d \cdot p_\infty .
\]
Let
\[
\pi_C \colon X_C \to C^{(d)},\qquad \pi_{C'} \colon X_{C'} \to C'^{(d)}
\]
be the blowups at $d \cdot p_\infty$ and $d \cdot p'_\infty$ respectively, with
exceptional divisors $E_C, E_{C'}$ and generic points $\eta_C, \eta_{C'}$.  The
local rings at these generic points are the discrete valuation rings
\[
V_C := \mathcal{O}_{X_C,\eta_C} \subset K(C^{(d)}),\qquad
V_{C'} := \mathcal{O}_{X_{C'},\eta_{C'}} \subset K(C'^{(d)}).
\]

% #############################################################################
\section{The master diagram}\label{sec:master}
% #############################################################################

The torsor tower over $U^{(d)}$ embeds into the projective geometry of $C^{(d)}$
and $C'^{(d)}$, which in turn carry the canonical blowups at the points at
infinity.  Assembling all of this:
\[
\begin{tikzcd}[row sep=2.4em, column sep=2.2em]
\cP^{[d]_H} \ar[r,"H"] \ar[d]
  & (\cP_{G/H})^{(d)} \ar[d,"q"] \ar[r,hookrightarrow,"\iota'"]
  & C'^{(d)} \ar[dd,"f^{(d)}"]
  & X_{C'} \ar[l,"\pi_{C'}"'] \ar[dd,dashed,"f_X"] \\
\cP^{[d]_G} \ar[r,"H"'] \ar[dr,"G"',bend right=10]
  & (\cP_{G/H})^{[d]} \ar[d,"G/H"]
  & & \\
  & U^{(d)} \ar[r,hookrightarrow,"\iota"']
  & C^{(d)}
  & X_C \ar[l,"\pi_C"]
\end{tikzcd}
\]
Reading the diagram from left to right:
\begin{description}
  \item[Left block (columns 1--2):] the torsor tower of
        \Cref{prop:torsor_tower_diagram} over $U^{(d)}$ --- $H$-torsors
        horizontally, the $G/H$-torsor on the lower right, and the composite
        $G$-torsor $\cP^{[d]_G} \to U^{(d)}$.
  \item[Middle block (column 3):] the compactifications.  Here $\iota$ is the
        canonical open immersion $U^{(d)} \hookrightarrow C^{(d)}$, and $\iota'$
        is the open immersion $(\cP_{G/H})^{(d)} \hookrightarrow C'^{(d)}$
        coming from $\cP_{G/H} \subset C'$.  Restricting $f^{(d)}$ along
        $\iota'$ recovers the canonical map $(\cP_{G/H})^{(d)} \to U^{(d)}$,
        i.e.\ the composite of $q$ with the $G/H$-torsor.
  \item[Right block (column 4):] the blowups at the diagonal points at infinity.
        The dashed arrow $f_X \colon X_{C'} \dashrightarrow X_C$ is the rational
        lift of $f^{(d)}$, defined on the strict transform; on generic points of
        the exceptional divisors it sends $\eta_{C'} \mapsto \eta_C$, inducing
        an extension of DVRs $V_C \hookrightarrow V_{C'}$.
\end{description}

% #############################################################################
\section{Unramifiedness of \texorpdfstring{$\cP^{[d]_G}$}{P\^{}[d]\_G} at \texorpdfstring{$\eta_C$}{eta\_C}}\label{sec:main}
% #############################################################################

From now on assume:

\begin{theorem}\label{thm:main}
%Suppose
\begin{enumerate}
  \item[\textnormal{(H1)}] $(\cP_{G/H})^{[d]} \to U^{(d)}$ is unramified at
        $\eta_C$; and
  \item[\textnormal{(H2)}] $\cP^{[d]_H} \to (\cP_{G/H})^{(d)}$ is unramified at
        $\eta_{C'}$.
  \item[\textnormal{(H3)}] $G=\Z/p^r\Z$ is cyclic of prime power order.
\end{enumerate}
%Then $\cP^{[d]_G} \to U^{(d)}$ is unramified at $\eta_C$.
\end{theorem}

\medskip
\noindent\textbf{Fields and places.}  With
$K_C := \Frac V_C$, $K_{C'} := \Frac V_{C'}$,
$F := K\bigl((\cP_{G/H})^{[d]}\bigr)$, $E := K(\cP^{[d]_G})$,
$M := E \cap K_{C'}$, and
$N := \Gal(E \cdot K_{C'}/K_{C'}) \cong \Gal(E/M) \leq H$, set
$V_F := V_{C'}|_F$, choose a place $v_E$ of $E$ above $V_F$ and a place
$v_\star$ of $E \cdot K_{C'}$ above both $v_E$ and $V_{C'}$, and write
$v_M := v_\star|_M = v_E|_M = V_{C'}|_M$.  The proof takes place inside the
parallel lattices of fields and of places lying over them:
\[
\begin{tikzcd}[row sep=1.6em, column sep=1.4em]
 & E \cdot K_{C'} \arrow[dl, dash] \arrow[dr, dash, "N"] & \\
E \arrow[dr, dash, "N"'] & & K_{C'} \arrow[dl, dash] \\
 & M \arrow[d, dash, "H/N"] & \\
 & F \arrow[d, dash, "G/H"] & \\
 & K_C &
\end{tikzcd}
\qquad\qquad
\begin{tikzcd}[row sep=1.6em, column sep=1.4em]
 & v_\star \arrow[dl, dash] \arrow[dr, dash] & \\
v_E \arrow[dr, dash] & & V_{C'} \arrow[dl, dash] \\
 & v_M \arrow[d, dash] & \\
 & V_F \arrow[d, dash] & \\
 & V_C &
\end{tikzcd}
\]
(lines denote field inclusion / place lying-over, read upwards; in the left
diagram a label is the Galois group of the extension where it is Galois).  In
addition $E/F$ is Galois with group $H$ and $E/K_C$ Galois with group $G$.

%------------------------------------------------------------------------------
\subsection{The ramification of \texorpdfstring{$v_M/V_F$}{vM/VF} and
            redundancy of \texorpdfstring{(H3)}{(H3)}}
\label{subsec:vM-VF-ramification}
%------------------------------------------------------------------------------

The set-up above leaves room, a priori, for $F \subsetneq M$, in which case the
ramification of $v_E/V_F$ would split nontrivially between the steps
$v_E/v_M$ and $v_M/V_F$.  In our geometric situation the Galois lattice
collapses this room: $M$ and $F$ coincide, $v_M / V_F$ is the trivial extension,
and (H3) is not used.

\begin{prop}\label{prop:M-equals-F}
$M = F$ as subfields of $E \cdot K_{C'}$.
\end{prop}

\begin{proof}
Write $L := K(\cP^{\times d})$.  Connectedness of $\cP^{\times d} \to U^d$
(reduction (R1)) and the abelianness of $G$ make $L/K_C$ Galois with group
$\widetilde G := G^d \rtimes S_d$, the $S_d$-action permuting coordinates.
Identifying the three subfields of $L$ via
$\cP^{[d]_G} = (K_G \rtimes S_d) \backslash \cP^{\times d}$,
$(\cP_{G/H})^{(d)} = (H^d \rtimes S_d) \backslash \cP^{\times d}$,
$(\cP_{G/H})^{[d]} = ((K_G + H^d) \rtimes S_d) \backslash \cP^{\times d}$, one
has
\[
E = L^{K_G \rtimes S_d}, \qquad
K_{C'} = L^{H^d \rtimes S_d}, \qquad
F = L^{(K_G + H^d) \rtimes S_d}.
\]
Galois correspondence identifies $E \cap K_{C'}$ with the fixed field of the
subgroup of $\widetilde G$ generated by $K_G \rtimes S_d$ and
$H^d \rtimes S_d$.  Inside the abelian group $G^d$ the subgroup generated by
$K_G$ and $H^d$ is the sum $K_G + H^d$, and $S_d$ normalises both factors
(each of $K_G$ and $H^d$ is permutation-invariant), so
\[
\langle K_G \rtimes S_d,\; H^d \rtimes S_d \rangle
   \;=\; (K_G + H^d) \rtimes S_d
   \;=\; \Gal(L/F).
\]
Hence $\Gal(L/M) = \Gal(L/F)$, i.e.\ $M = F$.
\end{proof}

\begin{cor}\label{cor:vM-trivial}
$v_M = V_F$ and the extension of discrete valuation rings $v_M / V_F$ is the
trivial extension; in particular it is unramified.
\end{cor}

\begin{proof}
By \Cref{prop:M-equals-F}, $M = F$.  By construction $v_M$ is the common
restriction $v_\star|_M = v_E|_M = V_{C'}|_M$; under $M = F$ each of these
restrictions equals $V_F = V_{C'}|_F$.
\end{proof}


\printbibliography

\end{document}
